%% =====================================================================
%%  T-22-03  The Familiar Face
%%  -------------------------------------------------------------------
%%  One idea per file. Core is mandatory. Slots in canonical order.
%%  Max 700 words. No layout commands. No filler. New source material
%%  for this entry goes in sources/B4/T-22-03.tex -- never by
%%  rewriting this file (docs/RULES.md R0.1).
%%  =====================================================================
%% @kind: topic
%% @id: T-22-03
%% @title: The Familiar Face
%% @subtitle: Manufactured familiarity --- we like what we have seen before, and cannot say why
%% @chapter: c22
%% @order: 30
%% @status: stable
%% @sources: B4
%% @updated: 2026-09-19

\topic{T-22-03}{The Familiar Face}{Manufactured familiarity --- we like what we have seen before, and cannot say why}

\begin{core}
Mostly, we like what we already recognise --- and the liking arrives without a reason attached. Mere repeated exposure raises liking and compliance, even when the exposure is beneath notice; voters pick the name they have seen; faces flashed too fast to remember still raise their owners' influence later. The manufactured-familiarity move is deliberate exposure engineering: arrange to be seen, again and slightly, until the sight of you is itself a warm signal.
\end{core}
\begin{mechanism}
B4's laboratory layer: subjects shown faces too briefly to recall later still liked the frequently flashed faces more, and were more persuaded by them --- the exposure counter ran without the subject's knowledge or memory. The mechanism is processing fluency: the mind mistakes the ease of recognising a thing for evidence about the thing, and familiarity is ease. B4's field cases: the candidate who changed his name to the state's political tradition and won; the mirror-image photograph experiment --- you prefer your reversed face, your friends the true print, because each prefers the version they have seen most. The operator's version is scheduled exposure: the same seat, the same corridor, the same name said in the right rooms, the logo on the screen before the pitch. B4's boundary condition is the finding that matters most: mere contact is not enough and can backfire --- school desegregation studies showed contact under unequal or competitive conditions increasing prejudice. Familiarity compounds liking only when the contact is cooperative and pleasant; under friction it compounds the opposite.
\end{mechanism}
\begin{conditions}
Strongest where decisions are made on low information --- name recognition elections, shelf choices, first rounds of hiring, the unknown vendor. Weakens where evaluation is deep (familiarity cannot survive a bad audit), and it reverses under overexposure: the over-familiar becomes wallpaper, then irritation. The cooperative condition is strict --- familiarity built inside conflict builds contempt.
\end{conditions}
\begin{application}
  \item Engineer the corridor: choose recurring, low-stakes co-presence --- same meeting, same gym hour, same forum --- before you need anything.
  \item Keep exposures short and pleasant. The dose is frequency, not depth; depth can wait for the relationship that frequency earns.
  \item Attach familiarity to the thing being sold: repeat the name, the face, the phrase. Fluency in the stimulus reads as truth of the proposition (T-15-01).
  \item Time exposures before judgment moments: the familiar candidate, the familiar supplier, wins ties they have no other right to.
  \item As defence, audit your warmth: for each person you inexplicably favour, count the exposures. Fluency is not evidence, and you are fluent in almost everyone you see weekly.
\end{application}
\begin{feedback}
  \item Working: strangers treat you as more trustworthy than your record supports --- and keep citing the sense that they know you.
  \item Working: you win the ties. Ties are what exposure engineering exists for.
  \item Failing: the exposures are noticed and resented --- \emph{why are you always here?} The dose broke the camouflage.
  \item Failing: your warmth for someone evaporates after one bad audit. It was fluency all along; good --- now apply the test prospectively.
\end{feedback}
\begin{failure}
The engineered version has a hollow centre: familiarity built without the cooperative floor collapses at first friction, and the collapse is worse than never having started --- the exposure made the disappointment personal. Chronic operators also overexpose: presence past the point of usefulness converts the warm signal into furniture, and furniture is eventually resented for taking up the room.
\end{failure}
\begin{countermeasures}
  \item Run the exposure count on your own preferences: candidates, brands, vendors, colleagues. Ask what you know besides their frequency.
  \item Distrust the ties that exposure wins: when choosing between two unknowns, add a day of evidence before adding a point for recognition.
  \item Watch the pre-pitch exposure play --- the sponsor everywhere for a quarter before the ask. Schedule your own due diligence to land cold (T-23-03).
  \item In your own selling, keep the exposure honest: be seen in association with delivery, not just in association. Familiarity plus one audit trail is durable; familiarity alone is a soap bubble with a schedule.
\end{countermeasures}

\sources{B4}
\seesources
\seealso{T-22-01, T-22-02, T-15-01}
